\documentclass[conference]{IEEEtran}
\IEEEoverridecommandlockouts
\usepackage[utf8]{inputenc}
\usepackage[T5]{fontenc}
\usepackage{cite}
\usepackage{amsmath,amssymb,amsfonts}
\usepackage{graphicx}
\usepackage{textcomp}
\usepackage{xcolor}
\usepackage{array}
\usepackage{url}

\def\BibTeX{{\rm B\kern-.05em{\sc i\kern-.025em b}\kern-.08em
    T\kern-.1667em\lower.7ex\hbox{E}\kern-.125emX}}

\begin{document}

\title{GangnamSign: Hệ Thống Hỗ Trợ Trẻ Khiếm Thính Học Ngôn Ngữ Ký Hiệu}

\author{\IEEEauthorblockN{Trần Xuân Phong}
\IEEEauthorblockA{\textit{HCI Course Project}\\
Mã sinh viên: 23021657\\
Vietnam}
\and
\IEEEauthorblockN{Bùi Minh Quang}
\IEEEauthorblockA{\textit{HCI Course Project}\\
Mã sinh viên: 23021669\\
Vietnam}
\and
\IEEEauthorblockN{Phạm Minh Quân}
\IEEEauthorblockA{\textit{HCI Course Project}\\
Vietnam}
\and
\IEEEauthorblockN{Trần Văn Vinh}
\IEEEauthorblockA{\textit{HCI Course Project}\\
Mã sinh viên: 23021749\\
Vietnam}
}

\maketitle

\begin{abstract}
Báo cáo này trình bày khung thiết kế và triển khai cho GangnamSign, một hệ thống hỗ trợ trẻ khiếm thính học ngôn ngữ ký hiệu thông qua bài học theo chủ đề, flashcard, quiz, luyện ký hiệu bằng camera và phản hồi tức thời. Nội dung được tổng hợp từ \textit{GangnamSign.pptx}, \textit{GangnamSign.pdf}, báo cáo tham khảo \textit{NHÓM 4 BÁO CÁO.pdf}, notebook huấn luyện mô hình và mã nguồn hiện có trong repository. Báo cáo tập trung vào bối cảnh nghiên cứu, khoảng trống của các giải pháp hiện có, nguyên lý thiết kế tương tác, đặc tả yêu cầu, kiến trúc hệ thống, pipeline nhận diện ký hiệu và định hướng phát triển. Kết quả từ notebook cho thấy mô hình landmark BiLSTM 46 lớp đạt validation accuracy 89.90\% và top-3 accuracy 96.10\%, tạo nền tảng cho chức năng nhận diện ký hiệu trong trải nghiệm học tập.
\end{abstract}

\begin{IEEEkeywords}
ngôn ngữ ký hiệu, trẻ khiếm thính, học tập trực quan, HCI, MediaPipe, BiLSTM, WLASL, phản hồi tức thời
\end{IEEEkeywords}

\section{Giới Thiệu}
Ngôn ngữ ký hiệu là một phương tiện giao tiếp quan trọng của cộng đồng người khiếm thính. Với trẻ em, việc tiếp cận ngôn ngữ ký hiệu từ sớm giúp phát triển tư duy ngôn ngữ, tăng khả năng giao tiếp với gia đình và hỗ trợ hòa nhập xã hội. Theo World Health Organization, gần 2.5 tỷ người được dự báo sẽ có một mức độ suy giảm thính lực vào năm 2050; hiện có khoảng 430 triệu người cần phục hồi chức năng thính lực, trong đó có 34 triệu trẻ em \cite{who}. Vì vậy, các công cụ học tập trực quan, dễ tiếp cận và có phản hồi nhanh là cần thiết đối với nhóm người dùng này.

GangnamSign được xây dựng theo hướng kết hợp giáo dục trực quan và tương tác người máy. Thay vì chỉ cung cấp tài liệu tĩnh, hệ thống tổ chức bài học theo topic, dùng flashcard và quiz để củng cố kiến thức, theo dõi tiến độ học tập và tích hợp chức năng nhận diện ký hiệu từ camera. Cách tiếp cận này kế thừa khung nội dung của bộ slide GangnamSign và cấu trúc báo cáo ứng dụng SignLearn trong \textit{NHÓM 4 BÁO CÁO.pdf}, đồng thời điều chỉnh theo mã nguồn và notebook của dự án hiện tại.

\section{Nguồn Tư Liệu và Phạm Vi}
Bảng~\ref{tab:sources} tóm tắt các nguồn được sử dụng để xây dựng báo cáo. Ngoài các slide, báo cáo tham khảo, notebook và mã nguồn, báo cáo còn sử dụng tài liệu tổng hợp \textit{tong\_hop\_artifact\_bao\_cao.docx}. Tài liệu này cung cấp phần diễn giải chi tiết hơn về bối cảnh, hiện trạng, người dùng, mục tiêu hệ thống, thiết kế giao diện, tương tác, kiến trúc và xử lý dữ liệu.

\begin{table}[htbp]
\caption{Nguồn tư liệu dùng để tổng hợp báo cáo}
\begin{center}
\footnotesize
\begin{tabular}{|p{0.31\linewidth}|p{0.57\linewidth}|}
\hline
\textbf{Nguồn} & \textbf{Vai trò trong báo cáo} \\
\hline
\textit{GangnamSign.pptx}/\textit{GangnamSign.pdf} & Bối cảnh nghiên cứu, vấn đề hiện tại, giải pháp đã có, cơ sở khoa học, nguyên lý thiết kế tương tác và khung demo. \\
\hline
\textit{NHÓM 4 BÁO CÁO.pdf} & Khung báo cáo đầy đủ: mô tả bài toán, đối tượng sử dụng, hướng tiếp cận, đặc tả yêu cầu, tính năng, UI, huấn luyện mô hình, khó khăn và định hướng. \\
\hline
\textit{notebook87071c7039.ipynb} & Dữ liệu, pipeline landmark, mô hình BiLSTM, tập lớp, cách chia dữ liệu và kết quả validation. \\
\hline
\textit{tong\_hop\_artifact\_bao\_cao.docx} & Nội dung tổng hợp cho các mục bối cảnh, hiện trạng, đối tượng sử dụng, mục tiêu, phạm vi, giải pháp, nguyên lý HCI, giao diện, tương tác, kiến trúc và xử lý dữ liệu. \\
\hline
Repository hiện tại & Công nghệ thực tế: React/React Router, FastAPI, PostgreSQL, SQLAlchemy, MediaPipe Tasks, TensorFlow/Keras, PyTorch, Gemini grading và API CRUD. \\
\hline
\end{tabular}
\label{tab:sources}
\end{center}
\end{table}

\section{Bối Cảnh và Bài Toán}
\subsection{Bối cảnh nghiên cứu}
Ngôn ngữ ký hiệu là phương tiện giao tiếp quan trọng đối với cộng đồng người khiếm thính. Đối với trẻ khiếm thính, việc tiếp cận ngôn ngữ ký hiệu từ sớm có ảnh hưởng trực tiếp đến khả năng giao tiếp, học tập và hòa nhập xã hội. Tuy nhiên, quá trình học thường gặp khó khăn do thiếu tài liệu phù hợp, thiếu môi trường luyện tập thường xuyên và thiếu công cụ phản hồi tức thời.

Các hình thức học truyền thống như học trực tiếp với giáo viên, học qua sách hoặc xem video minh họa vẫn có giá trị, nhưng chưa đáp ứng tốt nhu cầu học tập của trẻ em. Trẻ thường có khả năng tập trung ngắn và dễ mất hứng thú nếu nội dung học đơn điệu hoặc chỉ mang tính quan sát một chiều. Việc chỉ xem video và bắt chước cũng khiến trẻ khó biết mình thực hiện đúng hay sai nếu không có người hướng dẫn.

Trong bối cảnh đó, một hệ thống học ngôn ngữ ký hiệu dựa trên nền tảng web có thể hỗ trợ quá trình học theo hướng trực quan và tương tác hơn. Hệ thống có thể kết hợp video minh họa, bài học theo chủ đề, câu hỏi luyện tập, flashcard, trò chơi và phản hồi nhanh để giúp người học tiếp thu hiệu quả hơn.

\subsection{Hiện trạng vấn đề}
Trẻ khiếm thính thường gặp khó khăn trong giao tiếp xã hội, tiếp cận giáo dục và phát triển ngôn ngữ nếu thiếu môi trường học phù hợp. Các tài liệu học ngôn ngữ ký hiệu hiện nay vẫn thường ở dạng sách, ảnh tĩnh, video rời rạc hoặc nội dung thiếu tương tác. Khi trẻ chỉ xem video rồi bắt chước, trẻ khó biết mình thực hiện đúng hay sai vì thiếu phản hồi tức thời.

Một hạn chế lớn của nhiều tài liệu học hiện nay là nội dung được trình bày theo dạng quan sát một chiều. Người học xem hình ảnh hoặc video ký hiệu rồi tự ghi nhớ, nhưng hệ thống không đánh giá được mức độ hiểu bài, không chỉ ra lỗi sai và không cung cấp phản hồi ngay sau mỗi thao tác. Điều này khiến quá trình học dễ bị thụ động, đặc biệt với trẻ em cần phản hồi trực quan và nhanh để duy trì sự chú ý.

Đặc điểm nhận thức của trẻ cũng làm bài toán trở nên khác với người học trưởng thành. Trẻ có thời gian tập trung ngắn, học tốt hơn qua hình ảnh và chuyển động, cần nhiệm vụ nhỏ, mục tiêu rõ ràng và phản hồi đơn giản. Do đó, hệ thống không chỉ cần nhận diện hoặc hiển thị ký hiệu, mà còn phải thiết kế trải nghiệm học tập phù hợp với hành vi của trẻ.

\subsection{Đối tượng sử dụng}
Đối tượng chính của GangnamSign là trẻ khiếm thính khoảng 6--12 tuổi. Nhóm người dùng phụ gồm phụ huynh, giáo viên, người thân muốn học ký hiệu để giao tiếp với trẻ, và quản trị viên phụ trách tạo nội dung học tập. Việc xác định nhiều nhóm người dùng giúp hệ thống cân bằng giữa giao diện học tập đơn giản cho trẻ và công cụ quản trị đủ mạnh cho người lớn.

\subsection{Hướng tiếp cận}
Hệ thống hướng đến ba mục tiêu. Thứ nhất, cung cấp nội dung học theo chủ đề với hình ảnh, video, flashcard và quiz ngắn. Thứ hai, hỗ trợ luyện tập bằng camera để trẻ nhận phản hồi sau khi thực hiện ký hiệu. Thứ ba, cung cấp dashboard quản trị để tạo topic, lesson, quiz, theo dõi tiến độ và cập nhật nội dung.

\subsection{Mục tiêu hệ thống}
Mục tiêu chính của GangnamSign là xây dựng một nền tảng web hỗ trợ trẻ khiếm thính học ngôn ngữ ký hiệu theo hướng trực quan, ngắn gọn và có tương tác. Về mặt học tập, hệ thống giúp trẻ nhận biết, ghi nhớ và sử dụng các ký hiệu cơ bản trong giao tiếp. Nội dung được tổ chức theo từng chủ đề rõ ràng, mỗi bài học tập trung vào một đơn vị kiến thức nhỏ để tránh quá tải nhận thức.

Về mặt tương tác, hệ thống cần cung cấp phản hồi nhanh sau mỗi hoạt động. Khi trẻ trả lời câu hỏi hoặc hoàn thành bài luyện tập, hệ thống hiển thị kết quả đúng/sai, đáp án hoặc gợi ý phù hợp. Cơ chế này giúp trẻ nhận biết lỗi sai, củng cố trí nhớ và tăng động lực học tập. Về mặt quản lý, hệ thống hỗ trợ quản trị viên thêm, sửa, xóa và cập nhật topic, lesson, video ký hiệu, flashcard và câu hỏi luyện tập.

\subsection{Phạm vi báo cáo}
Báo cáo tập trung vào quá trình phân tích, thiết kế và xây dựng hệ thống web hỗ trợ trẻ khiếm thính học ngôn ngữ ký hiệu. Phạm vi chức năng gồm học theo chủ đề, xem nội dung ký hiệu bằng hình ảnh/video, luyện tập thông qua quiz, flashcard, theo dõi kết quả học tập và quản trị nội dung. Phạm vi kỹ thuật gồm frontend, backend, cơ sở dữ liệu và các thành phần xử lý dữ liệu hoặc mô hình đã triển khai/thử nghiệm.

Báo cáo không đi sâu vào chẩn đoán y học, đánh giá lâm sàng hoặc thay thế chương trình can thiệp chuyên môn. Hệ thống được xem là công cụ hỗ trợ học tập và luyện tập, không thay thế vai trò của giáo viên hoặc chuyên gia giáo dục đặc biệt.

\section{Giải Pháp Hiện Có và Khoảng Trống}
Các giải pháp truyền thống và công nghệ đều có giá trị, nhưng chưa giải quyết đầy đủ nhu cầu của trẻ em khiếm thính. Học trực tiếp với giáo viên có hiệu quả cao vì người học có thể được quan sát, chỉnh sửa động tác và giải thích ý nghĩa ký hiệu. Tuy nhiên, phương pháp này phụ thuộc vào thời gian, địa điểm, chi phí và số lượng giáo viên có chuyên môn.

Sách, tranh ảnh và tài liệu minh họa giúp người học tra cứu ký hiệu cơ bản, nhưng ngôn ngữ ký hiệu là dạng giao tiếp có chuyển động, biểu cảm và hướng tay nên hình ảnh tĩnh thường không truyền tải đầy đủ. Video hướng dẫn thể hiện chuyển động rõ hơn, nhưng phần lớn vẫn là học một chiều: hệ thống không biết người học đã làm đúng hay chưa và không phản hồi sau mỗi lần luyện tập. Bảng~\ref{tab:solutions} tổng hợp các nhóm giải pháp chính.

\begin{table}[htbp]
\caption{So sánh các nhóm giải pháp hiện có}
\begin{center}
\footnotesize
\begin{tabular}{|p{0.22\linewidth}|p{0.31\linewidth}|p{0.31\linewidth}|}
\hline
\textbf{Giải pháp} & \textbf{Ưu điểm} & \textbf{Hạn chế} \\
\hline
Giáo viên chuyên môn & Tương tác trực tiếp, phản hồi ngay, điều chỉnh theo từng trẻ & Số lượng giáo viên hạn chế, chi phí cao, phụ thuộc địa điểm và thời gian \\
\hline
Sách/từ điển ký hiệu & Chi phí thấp, dễ tiếp cận, có thể dùng ngoại tuyến & Ảnh tĩnh khó mô tả chuyển động tay, không có phản hồi, khó duy trì hứng thú \\
\hline
Cộng đồng người khiếm thính & Môi trường thực hành tự nhiên, hỗ trợ giao tiếp thật & Không phải lúc nào cũng có sẵn, khó duy trì lịch học đều đặn \\
\hline
Ứng dụng học ASL như Lingvano, The ASL App, HandSpeak, HandTalk & Có video, từ vựng phong phú, dễ tiếp cận trên thiết bị cá nhân & Nhiều sản phẩm hướng tới người lớn, ít cơ chế học dành riêng cho trẻ, hạn chế theo dõi tiến độ hoặc luyện tập có phản hồi \\
\hline
Nghiên cứu AI nhận diện ký hiệu & Có khả năng nhận diện cử chỉ và hỗ trợ đánh giá động tác & Thường dừng ở nguyên mẫu hoặc mô hình nhận diện, chưa gắn chặt với luồng học tập cho trẻ \\
\hline
\end{tabular}
\label{tab:solutions}
\end{center}
\end{table}

Khoảng trống chính là thiếu một nền tảng học ngôn ngữ ký hiệu được thiết kế chuyên biệt cho trẻ em, kết hợp nội dung trực quan, nhiệm vụ ngắn, game hóa, phản hồi tức thời và theo dõi tiến trình. Nhiều nền tảng có thể cung cấp từ điển hoặc video minh họa, nhưng chưa tổ chức thành bài học ngắn theo chủ đề, chưa có quiz/flashcard đủ rõ ràng hoặc chưa có cơ chế phản hồi giúp trẻ biết mình hiểu đúng hay sai. GangnamSign tập trung vào khoảng trống này thay vì chỉ làm một bộ từ điển hoặc một mô hình nhận diện đơn lẻ.

\section{Cơ Sở Khoa Học và Nguyên Lý HCI}
\subsection{Visual Learning Theory}
Visual Learning Theory cho rằng người học tiếp nhận và ghi nhớ thông tin hiệu quả hơn khi nội dung được trình bày qua hình ảnh, biểu tượng, chuyển động và các yếu tố trực quan. Đối với trẻ khiếm thính, kênh thị giác càng quan trọng vì quá trình tiếp nhận thông tin không phụ thuộc chủ yếu vào âm thanh. Do đó, hệ thống cần ưu tiên video ký hiệu, hình ảnh minh họa, biểu tượng dễ nhận biết và bố cục giao diện đơn giản.

Ngôn ngữ ký hiệu là một dạng ngôn ngữ có tính thị giác cao. Ý nghĩa của ký hiệu không chỉ nằm ở hình dạng bàn tay, mà còn phụ thuộc vào chuyển động, hướng tay, vị trí cơ thể và biểu cảm khuôn mặt. Nếu chỉ sử dụng văn bản hoặc hình ảnh tĩnh, người học có thể không hiểu đầy đủ cách thực hiện ký hiệu.

\subsection{Multimedia Learning Theory}
Theo Multimedia Learning Theory, người học tiếp thu tốt hơn khi thông tin được trình bày bằng sự kết hợp giữa hình ảnh, video, văn bản ngắn và hoạt động tương tác \cite{mayer}. Trong GangnamSign, nguyên lý này được áp dụng thông qua bài học có video ký hiệu, flashcard có hình ảnh hoặc từ khóa, quiz có lựa chọn trực quan và phản hồi sau mỗi câu trả lời.

Tuy nhiên, việc sử dụng đa phương tiện cần được kiểm soát hợp lý. Nếu một màn hình hiển thị quá nhiều video, chữ hoặc thành phần tương tác, trẻ có thể bị phân tán sự chú ý. Vì vậy, mỗi màn hình học nên tập trung vào một nội dung chính như một ký hiệu, một từ vựng hoặc một câu hỏi luyện tập.

\subsection{Cognitive Load Theory}
Cognitive Load Theory cho rằng bộ nhớ làm việc của con người có giới hạn \cite{sweller}. Với trẻ em, nguyên lý này càng quan trọng: mỗi màn hình nên tập trung vào một từ, một ký hiệu hoặc một nhiệm vụ. Trong học ngôn ngữ ký hiệu, trẻ cần đồng thời quan sát chuyển động, ghi nhớ hình dạng ký hiệu, hiểu ý nghĩa và thao tác với hệ thống. Nếu bài học chứa quá nhiều video, đoạn chữ dài hoặc nút chức năng không cần thiết, trẻ dễ mất tập trung và khó xác định nhiệm vụ chính.

\subsection{Feedback Learning}
Feedback Learning nhấn mạnh vai trò của phản hồi trong quá trình học tập. Người học cần biết hành động của mình đúng hay sai để điều chỉnh cách hiểu và củng cố kiến thức. Với trẻ khiếm thính, phản hồi cần trực quan, ngắn gọn và xuất hiện ngay sau thao tác. Nếu trẻ trả lời sai, hệ thống cần hiển thị đáp án đúng hoặc gợi ý. Nếu trẻ làm đúng, hệ thống có thể dùng điểm, sao, huy hiệu hoặc hiệu ứng xác nhận để tăng động lực.

\subsection{Heuristic thiết kế}
Các nguyên lý HCI được chọn gồm User-Centered Design, Visibility of System Status, Immediate Feedback và Consistency and Standards \cite{nielsen}. User-Centered Design yêu cầu hệ thống được thiết kế dựa trên đặc điểm, nhu cầu và giới hạn của trẻ khiếm thính. Visibility of System Status yêu cầu hệ thống cho người dùng biết mình đang ở đâu, câu trả lời đã được ghi nhận hay chưa, bài học đã hoàn thành hay chưa và bước tiếp theo là gì. Consistency and Standards yêu cầu các nút, biểu tượng, màu sắc và phản hồi có cùng ý nghĩa phải được dùng nhất quán trên toàn hệ thống.

Trong GangnamSign, các nguyên lý này được chuyển hóa thành: hiển thị tiến độ topic, giữ bố cục nhất quán giữa lesson/quiz/flashcard, dùng trạng thái rõ ràng khi camera đang ghi nhận, và đưa phản hồi đúng/sai ngay sau mỗi câu hỏi hoặc lần luyện ký hiệu.

\section{Đặc Tả Yêu Cầu}
\subsection{Yêu cầu chức năng}
Dựa trên tài liệu tổng hợp artifact và chức năng hiện có của GangnamSign, hệ thống cần các nhóm chức năng trong Bảng~\ref{tab:reqs}. Các chức năng này được chia theo vai trò người học, phụ huynh/giáo viên hỗ trợ và quản trị viên.

\begin{table}[htbp]
\caption{Yêu cầu chức năng chính}
\begin{center}
\footnotesize
\begin{tabular}{|p{0.29\linewidth}|p{0.55\linewidth}|}
\hline
\textbf{Nhóm chức năng} & \textbf{Mô tả} \\
\hline
Tài khoản và hồ sơ & Đăng ký, đăng nhập, tạo hồ sơ học sinh, bảo vệ API bằng xác thực người dùng. \\
\hline
Học theo topic & Hiển thị danh sách topic, bài học trong từng topic, trạng thái hoàn thành và tiến độ. \\
\hline
Flashcard và quiz & Cung cấp flashcard theo chủ đề, câu hỏi trắc nghiệm, câu hỏi theo video, đáp án đúng/sai và gợi ý ngắn. \\
\hline
Luyện ký hiệu bằng camera & Thu nhận ảnh hoặc landmark từ camera, gửi về backend để nhận diện ký hiệu hoặc chấm mức độ đúng. \\
\hline
Dịch ký hiệu và text-to-sign & Hỗ trợ sign-to-text cho chữ cái/từ vựng và text-to-sign thông qua SignWriting/pose/video khi có dữ liệu tương ứng. \\
\hline
Tiến độ và phần thưởng & Lưu bài đã hoàn thành, số câu đúng, XP, sao, streak, badge và lịch sử làm quiz. \\
\hline
Quản trị nội dung & Admin CRUD topic, lesson, exercise, quiz question, flashcard; xem báo cáo học sinh và thống kê bài học. \\
\hline
\end{tabular}
\label{tab:reqs}
\end{center}
\end{table}

\subsection{Yêu cầu phi chức năng}
Giao diện phải trực quan, thân thiện với trẻ, ưu tiên biểu tượng và màu sắc rõ ràng. Hệ thống cần phản hồi nhanh, tránh gián đoạn trong các luồng học chính, và bảo vệ dữ liệu cá nhân của người dùng. Với chức năng camera, hệ thống cần xử lý tình huống không có quyền truy cập camera, không phát hiện bàn tay hoặc ký hiệu nằm ngoài bộ nhãn.

\begin{table}[htbp]
\caption{Yêu cầu phi chức năng chính}
\begin{center}
\footnotesize
\begin{tabular}{|p{0.28\linewidth}|p{0.56\linewidth}|}
\hline
\textbf{Yêu cầu} & \textbf{Mô tả} \\
\hline
Dễ sử dụng & Giao diện đơn giản, nút lớn, nhãn ngắn, ưu tiên hình ảnh/video và hạn chế văn bản dài. \\
\hline
Phản hồi nhanh & Trả kết quả đúng/sai, trạng thái loading, lỗi dữ liệu và kết quả nhận diện trong thời gian đủ ngắn để không làm gián đoạn luồng học. \\
\hline
Nhất quán & Các màn hình lesson, quiz, flashcard, result và admin dùng cấu trúc điều hướng, màu trạng thái và kiểu phản hồi thống nhất. \\
\hline
An toàn dữ liệu & Kiểm soát truy cập bằng tài khoản, phân quyền và bảo vệ dữ liệu hồ sơ, tiến độ, lịch sử học và dữ liệu camera. \\
\hline
Dễ mở rộng & Nội dung học, câu hỏi, flashcard và model nhận diện có thể được bổ sung hoặc thay thế mà không cần thay đổi toàn bộ hệ thống. \\
\hline
\end{tabular}
\label{tab:nonfunctional}
\end{center}
\end{table}

\section{Thiết Kế Trải Nghiệm Người Dùng}
\subsection{Luồng học chính}
Luồng học được thiết kế theo chu trình ngắn: người học chọn topic, mở lesson, xem ký hiệu, luyện bằng flashcard hoặc quiz, nhận phản hồi, rồi cập nhật tiến độ. Cấu trúc này phù hợp với trẻ vì mỗi bước có mục tiêu rõ và thời lượng ngắn.

\subsection{Flashcard và quiz}
Flashcard giúp trẻ liên kết giữa từ vựng, hình ảnh và ký hiệu. Quiz dùng lựa chọn đúng/sai hoặc chọn đáp án theo video, giúp kiểm tra khả năng nhận diện ký hiệu. So với việc chỉ xem video, quiz tạo hành động chủ động và cho phép hệ thống ghi nhận kết quả học tập.

\subsection{Luyện ký hiệu bằng camera}
Khi luyện ký hiệu, người học đưa tay trước camera. Frontend kiểm tra sự xuất hiện của bàn tay, trích xuất keypoint hoặc gửi ảnh đã crop, sau đó backend trả về nhãn dự đoán, độ tin cậy và top-k nhãn gần nhất. Trong trường hợp mô hình cục bộ chưa đủ tin cậy, hệ thống có thể dùng Gemini để nhận diện hoặc chấm ký hiệu theo nhãn kỳ vọng.

\subsection{Thiết kế giao diện}
Giao diện nên sử dụng màu sắc tươi, tương phản rõ, font chữ dễ đọc và bố cục nhất quán. Các trạng thái học tập như progress, XP, sao và badge cần được hiển thị trực quan nhưng không làm rối màn hình. Với trẻ em, mỗi màn hình cần hạn chế số lựa chọn đồng thời để giảm tải nhận thức.

Bảng~\ref{tab:screens} tóm tắt các màn hình chính. Các màn hình học tập được tổ chức theo luồng rõ ràng, giảm thông tin thừa và ưu tiên phản hồi thị giác để phù hợp với người dùng chính.

\begin{table}[htbp]
\caption{Các màn hình chính của hệ thống}
\begin{center}
\footnotesize
\begin{tabular}{|p{0.26\linewidth}|p{0.30\linewidth}|p{0.28\linewidth}|}
\hline
\textbf{Màn hình} & \textbf{Người dùng} & \textbf{Mục đích} \\
\hline
Đăng nhập/đăng ký & Người học, phụ huynh, quản trị viên & Truy cập tài khoản và dữ liệu cá nhân. \\
\hline
Trang chủ học tập & Người học & Điều hướng đến topic, practice, translate và profile. \\
\hline
Danh sách topic & Người học & Chọn chủ đề học và xem tiến độ tổng quan. \\
\hline
Lesson/path & Người học & Xem bài học, ký hiệu, mô tả và trạng thái hoàn thành. \\
\hline
Quiz/flashcard & Người học & Luyện tập, kiểm tra ghi nhớ và nhận phản hồi. \\
\hline
Practice/translate & Người học & Luyện ký hiệu bằng camera hoặc chuyển đổi văn bản/ký hiệu. \\
\hline
Profile/progress & Người học, phụ huynh, giáo viên & Theo dõi kết quả, XP, badge và lịch sử học. \\
\hline
Admin & Quản trị viên & Quản lý topic, lesson, quiz, nội dung và báo cáo. \\
\hline
\end{tabular}
\label{tab:screens}
\end{center}
\end{table}

\subsection{Thiết kế tương tác}
Thiết kế tương tác tập trung vào cách người dùng thao tác với hệ thống và cách hệ thống phản hồi. Trong luồng học bài, người học bắt đầu bằng cách chọn một chủ đề, sau đó chọn bài học cụ thể. Hệ thống hiển thị video hoặc hình ảnh ký hiệu ở vị trí trung tâm, kèm tên ký hiệu và mô tả ngắn. Người học có thể phát lại nội dung, chuyển bài hoặc quay lại danh sách bài học.

Trong luồng quiz, tương tác chính là quan sát câu hỏi, chọn đáp án và nhận phản hồi. Sau khi người học chọn đáp án, hệ thống cần phản hồi ngay bằng trạng thái đúng/sai. Nếu đáp án sai, hệ thống hiển thị đáp án đúng hoặc gợi ý ngắn. Trong luồng flashcard, người học xem mặt trước của thẻ, lật thẻ để xem đáp án hoặc giải thích, rồi chuyển sang thẻ tiếp theo. Các tương tác này được giữ đơn giản để trẻ không phải xử lý nhiều thông tin cùng lúc.

\begin{table}[htbp]
\caption{Tương tác chính và phản hồi của hệ thống}
\begin{center}
\footnotesize
\begin{tabular}{|p{0.22\linewidth}|p{0.31\linewidth}|p{0.31\linewidth}|}
\hline
\textbf{Ngữ cảnh} & \textbf{Hành động} & \textbf{Phản hồi} \\
\hline
Học bài & Chọn chủ đề & Hiển thị danh sách bài học thuộc chủ đề. \\
\hline
Học bài & Chọn bài học & Hiển thị ký hiệu, hình/video và mô tả ngắn. \\
\hline
Quiz & Chọn đáp án & Hiển thị đúng/sai và đáp án đúng nếu cần. \\
\hline
Quiz & Hoàn thành quiz & Hiển thị kết quả tổng hợp và cập nhật tiến độ. \\
\hline
Flashcard & Lật thẻ & Hiển thị đáp án hoặc nội dung giải thích. \\
\hline
Camera & Thực hiện ký hiệu & Hiển thị nhãn dự đoán, confidence hoặc gợi ý lỗi. \\
\hline
Quản trị & Thêm/sửa nội dung & Lưu dữ liệu và hiển thị thông báo kết quả. \\
\hline
\end{tabular}
\label{tab:interactions}
\end{center}
\end{table}

\section{Kiến Trúc Hệ Thống}
\subsection{Tổng quan}
GangnamSign sử dụng kiến trúc client--server. Frontend đảm nhiệm giao diện học tập, camera, trích xuất keypoint và hiển thị phản hồi. Backend quản lý dữ liệu, xác thực, CRUD nội dung, tiến độ học tập và inference mô hình. Cơ sở dữ liệu lưu topic, lesson, exercise, quiz, hồ sơ học sinh, tiến độ và báo cáo.

\subsection{Frontend}
Frontend là thành phần chịu trách nhiệm hiển thị giao diện người dùng và xử lý các tương tác trực tiếp. Đây là phần đặc biệt quan trọng vì người dùng chính là trẻ khiếm thính, nhóm cần giao diện trực quan, thao tác đơn giản và phản hồi rõ ràng trong quá trình học.

Frontend hiện tại được xây dựng bằng React, React Router, TypeScript, Tailwind CSS, shadcn/Radix UI và MediaPipe Tasks JS. Các màn hình chính gồm trang học theo topic, lesson path, quiz, flashcard practice, translate/practice bằng camera, profile và admin dashboard. Khi người dùng thực hiện thao tác, frontend gửi request đến backend, nhận dữ liệu phản hồi và cập nhật giao diện tương ứng.

Một yêu cầu quan trọng của frontend là xử lý tốt trạng thái hệ thống. Khi dữ liệu đang tải, giao diện cần hiển thị trạng thái chờ. Khi người học chọn đáp án, hệ thống cần hiển thị phản hồi đúng/sai ngay. Khi xảy ra lỗi như không tải được nội dung, mất kết nối hoặc không phát hiện bàn tay, frontend cần thông báo ngắn và gợi ý thao tác tiếp theo. Trình duyệt cũng thực hiện một phần xử lý thị giác như phát hiện/crop bàn tay và trích xuất MediaPipe keypoints trước khi gửi dữ liệu về API.

\subsection{Backend}
Backend là thành phần xử lý nghiệp vụ, cung cấp API cho frontend và làm việc trực tiếp với cơ sở dữ liệu. Backend đảm bảo các chức năng như đăng nhập, lấy bài học, làm quiz, lưu tiến độ, nhận diện ký hiệu và quản lý nội dung được thực hiện đúng.

Backend hiện tại dùng FastAPI, SQLAlchemy async, Alembic, PostgreSQL, JWT authentication, TensorFlow/Keras, PyTorch và HTTP client cho dịch vụ AI ngoài. Các router chính gồm authentication, profiles, topics, lessons, questions, progress, reports và translate. Backend kiểm tra dữ liệu đầu vào, xử lý logic nghiệp vụ, truy xuất hoặc cập nhật cơ sở dữ liệu, sau đó trả kết quả về frontend theo định dạng thống nhất.

Một chức năng quan trọng của backend là xác thực và phân quyền. Người học chỉ được truy cập các chức năng học tập và tiến độ của hồ sơ thuộc tài khoản của mình, trong khi quản trị viên có quyền quản lý topic, lesson, quiz và nội dung học tập. Backend cũng cần chuẩn hóa xử lý lỗi như sai thông tin đăng nhập, thiếu quyền truy cập, không tìm thấy dữ liệu, dữ liệu đầu vào không hợp lệ hoặc lỗi kết nối cơ sở dữ liệu.

\subsection{Cơ sở dữ liệu}
Cơ sở dữ liệu lưu trữ toàn bộ thông tin cần thiết cho quá trình vận hành hệ thống, bao gồm tài khoản, hồ sơ học sinh, nội dung học tập, dữ liệu luyện tập, tiến độ và báo cáo. Trong artifact tổng hợp, phần database được mô tả theo hướng document database; tuy nhiên repository hiện tại triển khai bằng PostgreSQL kết hợp SQLAlchemy/Alembic. Vì vậy, báo cáo sử dụng PostgreSQL làm mô tả triển khai chính.

Các nhóm bảng chính gồm người dùng/hồ sơ, topic, lesson, exercise, quiz question, progress, attempt và badge. Một topic có nhiều lesson; một lesson có thể có câu hỏi quiz và dữ liệu luyện tập; một hồ sơ học sinh có nhiều bản ghi tiến độ và attempt. Các khóa như \texttt{user\_id}, \texttt{profile\_id}, \texttt{topic\_id} và \texttt{lesson\_id} được dùng để liên kết dữ liệu.

Cơ sở dữ liệu cần đảm bảo tính nhất quán trong các thao tác cập nhật. Khi quản trị viên thêm hoặc sửa bài học, dữ liệu cần được kiểm tra để tránh thiếu trường quan trọng như tên bài học, chủ đề liên kết, tài nguyên minh họa hoặc đáp án đúng. Khi người học hoàn thành bài học hoặc quiz, tiến độ cần được cập nhật đúng để hệ thống hiển thị kết quả chính xác.

\subsection{AI, Computer Vision và Model}
Thành phần AI/Computer Vision hỗ trợ nhận diện ký hiệu từ hình ảnh, camera hoặc dữ liệu landmark. MediaPipe được sử dụng để phát hiện bàn tay và trích xuất các điểm mốc. Mỗi bàn tay được biểu diễn bằng 21 hand landmarks; với pipeline landmark hiện tại, hệ thống còn kết hợp pose landmarks để bổ sung vị trí cơ thể.

Sau khi trích xuất landmarks, dữ liệu được chuyển thành vector hoặc chuỗi đặc trưng để đưa vào mô hình phân loại. Thành phần này có thể được tổ chức như một module riêng: frontend thu ảnh/frame hoặc keypoint, backend/module model chuẩn hóa đầu vào, gọi mô hình và trả về nhãn dự đoán cùng confidence/top-k. Cách tách riêng module nhận diện giúp hệ thống dễ bảo trì và dễ thay thế model khi có dữ liệu mới hoặc kiến trúc tốt hơn.

\begin{table}[htbp]
\caption{Luồng xử lý AI/Computer Vision}
\begin{center}
\footnotesize
\begin{tabular}{|p{0.14\linewidth}|p{0.28\linewidth}|p{0.40\linewidth}|}
\hline
\textbf{Bước} & \textbf{Thành phần} & \textbf{Mô tả} \\
\hline
1 & Frontend & Thu ảnh, frame hoặc keypoint từ camera người dùng. \\
\hline
2 & Backend/Model module & Nhận dữ liệu đầu vào và kiểm tra định dạng. \\
\hline
3 & MediaPipe/Preprocessing & Phát hiện bàn tay, trích xuất landmark và chuẩn hóa dữ liệu. \\
\hline
4 & Classification model & Dự đoán nhãn ký hiệu trong tập nhãn đã huấn luyện. \\
\hline
5 & Backend & Trả nhãn, confidence, top-k hoặc trạng thái lỗi. \\
\hline
6 & Frontend & Hiển thị kết quả nhận diện và phản hồi cho người học. \\
\hline
\end{tabular}
\label{tab:aiflow}
\end{center}
\end{table}

\begin{table}[htbp]
\caption{Công nghệ triển khai theo repository hiện tại}
\begin{center}
\footnotesize
\begin{tabular}{|p{0.28\linewidth}|p{0.56\linewidth}|}
\hline
\textbf{Lớp hệ thống} & \textbf{Công nghệ} \\
\hline
Frontend & React 19, React Router 7, TypeScript, Tailwind CSS, shadcn/Radix UI, lucide-react, MediaPipe Tasks JS, pose-viewer, Sutton SignWriting components. \\
\hline
Backend/API & FastAPI, Uvicorn, Pydantic, SQLAlchemy async, Alembic, JWT, passlib/bcrypt. \\
\hline
Dữ liệu & PostgreSQL; các bảng topic, lesson, exercise, quiz, profile, progress, attempt và badge. \\
\hline
ML/AI & TensorFlow/Keras cho mô hình alphabet và landmark BiLSTM; PyTorch Transformer cho WLASL word recognition; Gemini cho chấm ký hiệu khi cần. \\
\hline
\end{tabular}
\label{tab:tech}
\end{center}
\end{table}

\section{Huấn Luyện và Tích Hợp Mô Hình}
\subsection{Dữ liệu và lớp nhận diện}
Tài liệu tổng hợp artifact mô tả một pipeline nhận diện cử chỉ ở mức demo với dữ liệu hình ảnh được tổ chức theo thư mục nhãn. Mỗi thư mục con tương ứng với một nhãn ký hiệu, có thể đại diện cho chữ cái, con số hoặc một số từ/cử chỉ cơ bản. Artifact này ghi nhận tập dữ liệu thử nghiệm gồm 42 nhãn, huấn luyện trong 10 epoch và đạt validation accuracy khoảng 84.41\%.

Trong repository hiện tại, notebook huấn luyện mô hình landmark được mở rộng và kết hợp ba nhóm lớp: 10 từ WLASL (\textit{book}, \textit{finish}, \textit{go}, \textit{good}, \textit{help}, \textit{like}, \textit{mother}, \textit{what}, \textit{who}, \textit{yes}), 26 chữ cái ASL và 10 chữ số. Tổng số lớp là 46. Dữ liệu gồm video WLASL và các tập ảnh chữ cái/chữ số; ảnh tĩnh được lặp thành chuỗi 20 frame để phù hợp với đầu vào mô hình. WLASL là tập dữ liệu nhận diện ngôn ngữ ký hiệu Mỹ theo mức từ, được công bố trong WACV 2020 \cite{wlasl}.

\subsection{Trích xuất landmark}
Trước khi đưa dữ liệu vào mô hình, ảnh hoặc frame ký hiệu được tiền xử lý để trích xuất đặc trưng bàn tay. Thay vì huấn luyện trực tiếp trên toàn bộ ảnh thô, hệ thống sử dụng MediaPipe để phát hiện bàn tay và trích xuất các điểm mốc. Cách tiếp cận này giúp mô hình tập trung vào hình dạng, tư thế và chuyển động của tay thay vì nền ảnh hoặc chi tiết không liên quan.

Mỗi frame trong pipeline hiện tại được xử lý bằng MediaPipe để lấy landmark bàn tay và pose \cite{mediapipe}. Đặc trưng của một frame gồm bàn tay trái 21 điểm, bàn tay phải 21 điểm và pose 33 điểm; mỗi điểm có ba chiều $(x,y,z)$. Sau khi sắp xếp lại theo thứ tự bàn tay trái, bàn tay phải và pose, mỗi frame có 225 đặc trưng. Chuỗi đầu vào của mô hình có kích thước $20 \times 225$.

\subsection{Tiền xử lý}
Các chuỗi được chuẩn hóa theo từng frame. Bàn tay được neo tại cổ tay và chia theo khoảng cách landmark lớn nhất; pose được neo theo mốc cơ thể và chuẩn hóa theo độ rộng vai. Với video có nhiều hơn 20 frame, hệ thống lấy mẫu đều; với dữ liệu ngắn hơn, hệ thống lặp hoặc pad frame cuối. Cách tiền xử lý này giúp giảm ảnh hưởng của vị trí người học trong camera.

Quy trình tiền xử lý tổng quát gồm đọc dữ liệu từ thư mục nhãn hoặc video, gán nhãn theo tên lớp, phát hiện bàn tay bằng MediaPipe, trích xuất landmark, chuyển landmark thành vector/chuỗi đặc trưng và chuẩn hóa định dạng trước khi huấn luyện.

\subsection{Mô hình BiLSTM}
Mô hình được xây dựng theo hướng phân loại đa lớp: đầu vào là đặc trưng landmark, đầu ra là nhãn ký hiệu tương ứng. Artifact DOCX mô tả cấu hình thử nghiệm với Categorical Crossentropy, learning rate ban đầu 0.001, dropout 0.05, 10 epoch và chia dữ liệu 80\% training, 10\% validation, 10\% test.

Mô hình chính trong notebook hiện tại gồm Layer Normalization, hai tầng Bidirectional LSTM, Dropout, Dense, Batch Normalization và softmax 46 lớp. Tổng số tham số là 573,168. Dữ liệu sau trích xuất có 6,885 mẫu thuộc đủ 46 lớp; sau cân bằng còn 142 mẫu mỗi lớp, tạo 6,532 mẫu. Tập train có 5,225 mẫu và tập validation có 1,307 mẫu.

\begin{table}[htbp]
\caption{Tóm tắt kết quả huấn luyện từ notebook}
\begin{center}
\footnotesize
\begin{tabular}{|p{0.38\linewidth}|p{0.44\linewidth}|}
\hline
\textbf{Hạng mục} & \textbf{Giá trị} \\
\hline
Số lớp & 46 lớp: 10 từ WLASL, A--Z, 0--9 \\
\hline
Đầu vào mô hình & 20 frame $\times$ 225 đặc trưng landmark \\
\hline
Số mẫu trước cân bằng & 6,885 mẫu \\
\hline
Số mẫu sau cân bằng & 6,532 mẫu \\
\hline
Train/Validation & 5,225 / 1,307 mẫu \\
\hline
Kiến trúc & BiLSTM, Dense, BatchNorm, Dropout, Softmax \\
\hline
Validation accuracy & 89.90\% (1175/1307) \\
\hline
Top-3 accuracy & 96.10\% \\
\hline
Random baseline & 2.2\% \\
\hline
\end{tabular}
\label{tab:model}
\end{center}
\end{table}

\begin{table}[htbp]
\caption{Kết quả mô hình trong artifact DOCX}
\begin{center}
\footnotesize
\begin{tabular}{|p{0.38\linewidth}|p{0.44\linewidth}|}
\hline
\textbf{Chỉ số} & \textbf{Kết quả} \\
\hline
Số nhãn & 42 \\
\hline
Số epoch & 10 \\
\hline
Loss ban đầu & khoảng 2.65 \\
\hline
Loss cuối & khoảng 0.73 \\
\hline
Validation accuracy & khoảng 84.41\% \\
\hline
\end{tabular}
\label{tab:artifact-model}
\end{center}
\end{table}

\subsection{Tích hợp vào hệ thống}
Sau khi huấn luyện, mô hình được lưu lại và tích hợp vào hệ thống để phục vụ chức năng nhận diện ký hiệu. Người dùng cung cấp đầu vào thông qua camera hoặc hình ảnh; hệ thống xử lý đầu vào, trích xuất landmark bằng MediaPipe, chuẩn hóa đặc trưng rồi đưa vào mô hình để dự đoán nhãn.

Backend cung cấp các endpoint nhận diện ký hiệu: ảnh đơn cho alphabet CNN, keypoint clip cho WLASL Transformer 300 glosses, và landmark clip cho BiLSTM 46 lớp. Với endpoint landmark, frontend gửi dữ liệu dạng $(T,75,3)$ theo thứ tự pose, left hand và right hand; backend đổi thứ tự, chuẩn hóa, ép độ dài về 20 frame rồi chạy Keras inference. Kết quả trả về gồm nhãn dự đoán, confidence, top-k và trạng thái model loaded. Nếu mô hình chưa đủ ổn định, chức năng nhận diện nên được trình bày như tính năng thử nghiệm hỗ trợ luyện tập, chưa phải cơ chế đánh giá duy nhất.

\section{Demo và Luồng Kiểm Thử}
Demo hệ thống nên thể hiện ba luồng chính. Luồng học tập: người học chọn topic, mở lesson, xem ký hiệu, làm quiz và thấy tiến độ tăng. Luồng flashcard: người học chọn topic, lật thẻ, xem từ vựng/video và luyện ghi nhớ. Luồng camera: người học bật camera, thực hiện ký hiệu, hệ thống trích keypoint, dự đoán nhãn và hiển thị phản hồi.

Kiểm thử chức năng cần bao phủ các tình huống: đăng nhập/đăng ký, tạo hồ sơ học sinh, tải topic/lesson, hoàn thành bài học, lưu attempt, làm topic quiz, nhận diện ký hiệu khi có bàn tay, xử lý khi không có bàn tay, và thao tác CRUD admin. Kiểm thử HCI nên tập trung vào khả năng trẻ hiểu màn hình, thời gian hoàn thành bài học, mức độ nhầm lẫn khi dùng camera và độ rõ của phản hồi đúng/sai.

\section{Kết Quả Hiện Tại, Khó Khăn và Định Hướng}
\subsection{Kết quả hiện tại}
Hệ thống đã có nền tảng học theo topic, lesson, quiz, flashcard, progress, admin CRUD và API nhận diện ký hiệu. Các chức năng có thể demo gồm đăng ký/đăng nhập, tạo hồ sơ học sinh, xem danh sách topic, mở lộ trình bài học, xem lesson, làm quiz, xem kết quả, theo dõi tiến độ, luyện ký hiệu bằng camera, dùng màn hình translate/practice và quản trị topic/lesson/quiz.

Backend đã cung cấp các nhóm API cho auth, profile, topic, lesson, quiz, progress, report và translate. Dữ liệu seed ban đầu gồm 5 topic, 38 lesson, 24 câu hỏi quiz và 3 badge. Notebook đã huấn luyện được mô hình landmark BiLSTM có kết quả validation tốt trên tập dữ liệu cân bằng. Backend cũng có các dịch vụ nhận diện alphabet, WLASL word-level và chấm ký hiệu bằng Gemini.

\begin{table}[htbp]
\caption{Đánh giá theo yêu cầu ban đầu}
\begin{center}
\footnotesize
\begin{tabular}{|p{0.31\linewidth}|p{0.20\linewidth}|p{0.33\linewidth}|}
\hline
\textbf{Yêu cầu} & \textbf{Trạng thái} & \textbf{Ghi chú} \\
\hline
Đăng ký, đăng nhập, hồ sơ học sinh & Đã hoàn thành & Có frontend route và backend API tương ứng. \\
\hline
Học theo topic và lesson & Đã hoàn thành & Có danh sách topic, lộ trình bài học và dữ liệu seed. \\
\hline
Quiz và phản hồi đúng/sai & Đã hoàn thành & Có quiz theo lesson, topic quiz và result page. \\
\hline
Theo dõi tiến độ, XP, badge & Đã hoàn thành & Có progress API, attempt, XP, sao, streak và badge. \\
\hline
Luyện ký hiệu bằng camera & Hoàn thành một phần & Có camera/keypoint và API nhận diện; cần kiểm thử người dùng thật. \\
\hline
Text-to-sign/sign-to-text & Hoàn thành một phần & Có translate route và dịch vụ SignWriting/pose/model tùy dữ liệu. \\
\hline
Quản trị nội dung & Đã hoàn thành & Có admin overview, content, topics, lessons và quizzes. \\
\hline
Đánh giá định lượng người dùng & Chưa hoàn thiện & Chưa có usability testing và số liệu định lượng đầy đủ. \\
\hline
\end{tabular}
\label{tab:req-status}
\end{center}
\end{table}

\subsection{Khó khăn}
Ngôn ngữ ký hiệu là bài toán khó vì ký hiệu phụ thuộc không chỉ vào bàn tay mà còn vào chuyển động, vị trí cơ thể, biểu cảm khuôn mặt và ngữ cảnh. Dữ liệu chất lượng cao cho trẻ em và ký hiệu liên tục còn hạn chế. Dữ liệu camera thực tế có nhiều biến thiên về ánh sáng, góc quay, khoảng cách, tốc độ thực hiện ký hiệu và hình dạng bàn tay.

Về mô hình, kết quả validation trong notebook chưa đủ để khẳng định độ ổn định trong mọi điều kiện sử dụng thật. Một số ký hiệu có hình dạng bàn tay gần giống nhau dễ bị nhầm lẫn; các ký hiệu có chuyển động liên tục cần mô hình nắm bắt trình tự thời gian tốt hơn. Hệ thống cũng kết hợp nhiều pipeline như alphabet CNN, WLASL Transformer, landmark BiLSTM và Gemini, làm tăng độ phức tạp khi tích hợp, debug và đánh giá.

Về giao diện, người dùng chính là trẻ em nên mỗi màn hình phải rõ ràng, ít tải nhận thức và có phản hồi trực quan. Hệ thống chưa có kiểm thử khả dụng với trẻ khiếm thính, phụ huynh hoặc giáo viên nên chưa thể kết luận giao diện đã thật sự dễ dùng. Chức năng camera cũng cần hướng dẫn rõ hơn khi không phát hiện tay, ánh sáng yếu hoặc người học đặt tay ngoài khung hình.

\subsection{Định hướng phát triển}
Các hướng phát triển tiếp theo gồm mở rộng dữ liệu ký hiệu với nhiều từ vựng, nhiều người thực hiện và nhiều điều kiện quay khác nhau; bổ sung dữ liệu ký hiệu liên tục và các mẫu câu giao tiếp hằng ngày; cá nhân hóa lộ trình học dựa trên tiến độ, số lần sai, tốc độ học và độ khó của từng người học.

Hệ thống cũng cần thêm feedback chi tiết khi người học ký hiệu sai, ví dụ sai hình dạng bàn tay, sai vị trí tay, sai hướng chuyển động hoặc thiếu frame. Về mô hình, cần cải thiện nhận diện gesture bằng cách kết hợp hand landmark, pose, biểu cảm khuôn mặt và motion features; đánh giá model trên tập test độc lập, camera thực tế và nhiều nhóm người dùng. Về sản phẩm, cần hỗ trợ đa ngôn ngữ, mở rộng kho hình/video/pose/SignWriting, tối ưu độ trễ camera, bổ sung chế độ offline/cache và tăng cường bảo mật dữ liệu người học.

Quan trọng nhất, hệ thống cần được kiểm thử với người dùng thật, bao gồm trẻ khiếm thính, phụ huynh và giáo viên. Các chỉ số cần thu thập gồm thời gian hoàn thành bài học, tỷ lệ trả lời đúng, tần suất lỗi camera, mức độ cải thiện sau nhiều lần học và mức độ hài lòng của người dùng.

\section{Kết Luận}
GangnamSign hướng đến một hệ thống học ngôn ngữ ký hiệu trực quan, tương tác và phù hợp với trẻ khiếm thính. Dựa trên slide GangnamSign, báo cáo SignLearn, tài liệu tổng hợp artifact, notebook huấn luyện và mã nguồn hiện tại, báo cáo đã trình bày bối cảnh, khoảng trống giải pháp, nguyên lý HCI, đặc tả yêu cầu, thiết kế giao diện, kiến trúc client--server, pipeline nhận diện landmark và kết quả mô hình.

Trọng tâm của hệ thống không chỉ là nhận diện ký hiệu, mà là tạo một trải nghiệm học tập có phản hồi, theo dõi tiến độ và đủ linh hoạt để giáo viên/quản trị viên cập nhật nội dung. Các chức năng như đăng nhập, hồ sơ học sinh, học theo topic, làm quiz, theo dõi tiến độ, luyện ký hiệu bằng camera và quản trị nội dung đã có thể demo được. Tuy nhiên, hệ thống vẫn cần mở rộng dữ liệu, cải thiện độ ổn định của model, kiểm thử giao diện với người dùng thật và bổ sung đánh giá định lượng trước khi có thể kết luận về hiệu quả trong môi trường học tập thực tế.

\section*{Acknowledgment}
Nhóm tác giả cảm ơn giảng viên học phần Tương tác người máy và các bạn trong lớp đã góp ý trong quá trình phát triển đề tài.

\begin{thebibliography}{00}
\bibitem{who} World Health Organization, ``Deafness and hearing loss,'' Mar. 3, 2026. [Online]. Available: \url{https://www.who.int/news-room/fact-sheets/detail/deafness-and-hearing-loss}
\bibitem{mayer} R. E. Mayer, \textit{Multimedia Learning}, 2nd ed. Cambridge, U.K.: Cambridge University Press, 2009.
\bibitem{sweller} J. Sweller, ``Cognitive load during problem solving: Effects on learning,'' \textit{Cognitive Science}, vol. 12, no. 2, pp. 257--285, 1988.
\bibitem{nielsen} J. Nielsen, ``10 usability heuristics for user interface design,'' Nielsen Norman Group, Apr. 24, 1994, updated Jan. 30, 2024. [Online]. Available: \url{https://www.nngroup.com/articles/ten-usability-heuristics/}
\bibitem{wlasl} D. Li, C. Rodriguez, X. Yu, and H. Li, ``Word-level deep sign language recognition from video: A new large-scale dataset and methods comparison,'' in \textit{Proc. IEEE Winter Conf. Applications of Computer Vision (WACV)}, 2020, pp. 1459--1469.
\bibitem{mediapipe} C. Lugaresi \textit{et al.}, ``MediaPipe: A framework for building perception pipelines,'' arXiv:1906.08172, 2019.
\end{thebibliography}

\end{document}
